\documentclass[10pt,a4paper]{amsart}
\usepackage{amsmath,amssymb,amsthm,graphicx}
% Retained preamble for the Stacks Project chapter {CH} (shared project preamble).
% The chapter file pinned-original-source.tex only carries \input{{preamble}} on its first line,
% so no preamble text lives inside the chapter itself.  The Stacks Project uses one single shared
% preamble for every chapter: the snapshot stages it three times (sources/stacks-chow/preamble.tex,
% sources/stacks-etale-cohomology/preamble.tex and focused-seeds/AGP/sources/stacks-snapshot-20261011/
% preamble.tex) and those copies are identical in size.  The block below is copied mechanically,
% line for line, from sources/stacks-chow/preamble.tex lines 14-18 and 154-217: the Xy-pic setup and
% the native theorem declarations together with the mathematical macros.  Nothing was added, removed
% or redefined; the wrapper class and the amsmath/amssymb/amsthm/graphicx base packages are supplied
% by the assembler.
% For commutative diagrams we use Xy-pic
% Chapter: stacks-spaces-more-morphisms.tex -- More on Morphisms of Spaces.
\usepackage[all]{xy}

% We use 2cell for 2-commutative diagrams.
\xyoption{2cell}

% Theorem environments.
%
\theoremstyle{plain}
\newtheorem{theorem}[subsection]{Theorem}
\newtheorem{proposition}[subsection]{Proposition}
\newtheorem{lemma}[subsection]{Lemma}

\theoremstyle{definition}
\newtheorem{definition}[subsection]{Definition}
\newtheorem{example}[subsection]{Example}
\newtheorem{exercise}[subsection]{Exercise}
\newtheorem{situation}[subsection]{Situation}

\theoremstyle{remark}
\newtheorem{remark}[subsection]{Remark}
\newtheorem{remarks}[subsection]{Remarks}

\numberwithin{equation}{subsection}

% Macros
%
\def\lim{\mathop{\mathrm{lim}}\nolimits}
\def\colim{\mathop{\mathrm{colim}}\nolimits}
\def\Spec{\mathop{\mathrm{Spec}}}
\def\Hom{\mathop{\mathrm{Hom}}\nolimits}
\def\Ext{\mathop{\mathrm{Ext}}\nolimits}
\def\SheafHom{\mathop{\mathcal{H}\!\mathit{om}}\nolimits}
\def\SheafExt{\mathop{\mathcal{E}\!\mathit{xt}}\nolimits}
\def\Sch{\mathit{Sch}}
\def\Mor{\mathop{\mathrm{Mor}}\nolimits}
\def\Ob{\mathop{\mathrm{Ob}}\nolimits}
\def\Sh{\mathop{\mathit{Sh}}\nolimits}
\def\NL{\mathop{N\!L}\nolimits}
\def\CH{\mathop{\mathrm{CH}}\nolimits}
\def\proetale{{pro\text{-}\acute{e}tale}}
\def\etale{{\acute{e}tale}}
\def\QCoh{\mathit{QCoh}}
\def\Ker{\mathop{\mathrm{Ker}}}
\def\Im{\mathop{\mathrm{Im}}}
\def\Coker{\mathop{\mathrm{Coker}}}
\def\Coim{\mathop{\mathrm{Coim}}}

% Boxtimes
%
\DeclareMathSymbol{\boxtimes}{\mathbin}{AMSa}{"02}

%
% Macros for moduli stacks/spaces
%
\def\QCohstack{\mathcal{QC}\!\mathit{oh}}
\def\Cohstack{\mathcal{C}\!\mathit{oh}}
\def\Spacesstack{\mathcal{S}\!\mathit{paces}}
\def\Quotfunctor{\mathrm{Quot}}
\def\Hilbfunctor{\mathrm{Hilb}}
\def\Curvesstack{\mathcal{C}\!\mathit{urves}}
\def\Polarizedstack{\mathcal{P}\!\mathit{olarized}}
\def\Complexesstack{\mathcal{C}\!\mathit{omplexes}}
% \Pic is the operator that assigns to X its picard group, usage \Pic(X)
% \Picardstack_{X/B} denotes the Picard stack of X over B
% \Picardfunctor_{X/B} denotes the Picard functor of X over B
\def\Pic{\mathop{\mathrm{Pic}}\nolimits}
\def\Picardstack{\mathcal{P}\!\mathit{ic}}
\def\Picardfunctor{\mathrm{Pic}}
\def\Deformationcategory{\mathcal{D}\!\mathit{ef}}
\usepackage{cleveref}
\crefname{originalreference}{original reference}{original references}
\Crefname{originalreference}{Original reference}{Original references}
\makeatletter
\newlabel{decent-spaces-lemma-check-dimension-function-finite-cover}{{0ED9}{source}{}{sourcepointer-1}{}}
\newlabel{decent-spaces-lemma-check-dimension-function-finite-cover@cref}{{[originalreference][0][]0ED9}{[1][0][]source}{}{}{}}
\makeatother
\crefname{originalreference}{original reference}{original references}
\Crefname{originalreference}{Original reference}{Original references}
\makeatletter
\newlabel{decent-spaces-lemma-conditions-on-fibre-and-qf}{{0ACK}{source}{}{sourcepointer-2}{}}
\newlabel{decent-spaces-lemma-conditions-on-fibre-and-qf@cref}{{[originalreference][0][]0ACK}{[1][0][]source}{}{}{}}
\makeatother
\crefname{originalreference}{original reference}{original references}
\Crefname{originalreference}{Original reference}{Original references}
\makeatletter
\newlabel{decent-spaces-lemma-decent-space-closed-point}{{0AHB}{source}{}{sourcepointer-3}{}}
\newlabel{decent-spaces-lemma-decent-space-closed-point@cref}{{[originalreference][0][]0AHB}{[1][0][]source}{}{}{}}
\makeatother
\crefname{originalreference}{original reference}{original references}
\Crefname{originalreference}{Original reference}{Original references}
\makeatletter
\newlabel{decent-spaces-lemma-locally-Noetherian-decent-quasi-separated}{{0BB6}{source}{}{sourcepointer-4}{}}
\newlabel{decent-spaces-lemma-locally-Noetherian-decent-quasi-separated@cref}{{[originalreference][0][]0BB6}{[1][0][]source}{}{}{}}
\makeatother
\crefname{originalreference}{original reference}{original references}
\Crefname{originalreference}{Original reference}{Original references}
\makeatletter
\newlabel{decent-spaces-lemma-scheme-with-dimension-function}{{0ED5}{source}{}{sourcepointer-5}{}}
\newlabel{decent-spaces-lemma-scheme-with-dimension-function@cref}{{[originalreference][0][]0ED5}{[1][0][]source}{}{}{}}
\makeatother
\crefname{originallemma}{original lemma}{original lemmas}
\Crefname{originallemma}{Original lemma}{Original lemmas}
\makeatletter
\newlabel{lemma-construct-glueing}{{0EDM}{source}{}{sourcepointer-6}{}}
\newlabel{lemma-construct-glueing@cref}{{[originallemma][0][]0EDM}{[1][0][]source}{}{}{}}
\makeatother
\crefname{originalreference}{original reference}{original references}
\Crefname{originalreference}{Original reference}{Original references}
\makeatletter
\newlabel{properties-lemma-locally-Noetherian-closed-point}{{02IL}{source}{}{sourcepointer-7}{}}
\newlabel{properties-lemma-locally-Noetherian-closed-point@cref}{{[originalreference][0][]02IL}{[1][0][]source}{}{}{}}
\makeatother
\crefname{originalreference}{original reference}{original references}
\Crefname{originalreference}{Original reference}{Original references}
\makeatletter
\newlabel{spaces-limits-proposition-there-is-a-scheme-finite-over}{{09YC}{source}{}{sourcepointer-8}{}}
\newlabel{spaces-limits-proposition-there-is-a-scheme-finite-over@cref}{{[originalreference][0][]09YC}{[1][0][]source}{}{}{}}
\makeatother
\crefname{originalreference}{original reference}{original references}
\Crefname{originalreference}{Original reference}{Original references}
\makeatletter
\newlabel{spaces-morphisms-definition-dimension-fibre}{{04NM}{source}{}{sourcepointer-9}{}}
\newlabel{spaces-morphisms-definition-dimension-fibre@cref}{{[originalreference][0][]04NM}{[1][0][]source}{}{}{}}
\makeatother
\crefname{originalreference}{original reference}{original references}
\Crefname{originalreference}{Original reference}{Original references}
\makeatletter
\newlabel{spaces-morphisms-lemma-dimension-fibre-at-a-point-additive}{{04NR}{source}{}{sourcepointer-10}{}}
\newlabel{spaces-morphisms-lemma-dimension-fibre-at-a-point-additive@cref}{{[originalreference][0][]04NR}{[1][0][]source}{}{}{}}
\makeatother
\crefname{originalreference}{original reference}{original references}
\Crefname{originalreference}{Original reference}{Original references}
\makeatletter
\newlabel{spaces-morphisms-lemma-finite-proper}{{04NZ}{source}{}{sourcepointer-11}{}}
\newlabel{spaces-morphisms-lemma-finite-proper@cref}{{[originalreference][0][]04NZ}{[1][0][]source}{}{}{}}
\makeatother
\title{The transcendence-degree dimension function on a decent space over a universally catenary base}
\author{The Stacks Project Authors}
\date{stacks-project:spaces-more-morphisms.tex:local-snapshot-20261011; GNU FDL 1.2 or later}
\begin{document}
\maketitle
\noindent\textit{Editorial note.} Counted span is the author's own proof environment, source lines 7365-7457, immediately following the statement at 7348-7363.\par\medskip
\noindent\textit{Editorial note.} The statement contains the author's typo 'transcendence degreeof x/f(x)' (missing space); it is retained verbatim.\par\medskip
\noindent\textit{Editorial note.} Ten cross-chapter pointers are recorded in assembly.external\_references; all are navigation pointers and the chain-lifting argument itself is inside the span.\par\medskip
\noindent\textit{Allowed standard prerequisites (editorial).}\begin{itemize}
\item Dimension functions and their additivity along fibres. Morphisms of Spaces, definition spaces-morphisms-definition-dimension-fibre and lemma spaces-morphisms-lemma-dimension-fibre-at-a-point-additive, cited in the statement and proof. Defines delta\_X as delta plus transcendence degree and supplies additivity along specialisations inside fibres.
\item Catenarity of the base transfers to the total space. Decent Spaces lemmas decent-spaces-lemma-locally-Noetherian-decent-quasi-separated and decent-spaces-lemma-check-dimension-function-finite-cover, cited inside the proof. Reduces catenarity questions to a finite cover and lets the argument work on decent spaces rather than schemes.
\item Existence of a scheme finite over the space. Limits of Spaces, proposition spaces-limits-proposition-there-is-a-scheme-finite-over, cited inside the proof. Produces the finite morphism pi' : Y' -> X used to lift the two chains of specialisations to a common object.
\item Constructing the gluing used in the reduction. chapter lemma lemma-construct-glueing, cited inside the proof. Builds the finite Y' with prescribed behaviour away from the fibre over the chosen point.
\end{itemize}\medskip
\ifdefined\hypertarget\hypertarget{sourcepointer-1}{}\fi
\noindent\textit{Source reference.} Original reference 0ED9: focused-seeds/geometry/S9-spaces-univ-catenary-dim-func/ref-tags.txt, source lines 4--4. Navigation pointer only: the counted proof cites this statement by \textbackslash{}ref as a standing input; the deduction the counted proof has to make is carried out inside the counted span, and no part of that deduction is delegated to this pointer.\par\medskip
\ifdefined\hypertarget\hypertarget{sourcepointer-2}{}\fi
\noindent\textit{Source reference.} Original reference 0ACK: focused-seeds/geometry/S9-spaces-univ-catenary-dim-func/ref-tags.txt, source lines 5--5. Navigation pointer only: the counted proof cites this statement by \textbackslash{}ref as a standing input; the deduction the counted proof has to make is carried out inside the counted span, and no part of that deduction is delegated to this pointer.\par\medskip
\ifdefined\hypertarget\hypertarget{sourcepointer-3}{}\fi
\noindent\textit{Source reference.} Original reference 0AHB: focused-seeds/geometry/S9-spaces-univ-catenary-dim-func/ref-tags.txt, source lines 6--6. Navigation pointer only: the counted proof cites this statement by \textbackslash{}ref as a standing input; the deduction the counted proof has to make is carried out inside the counted span, and no part of that deduction is delegated to this pointer.\par\medskip
\ifdefined\hypertarget\hypertarget{sourcepointer-4}{}\fi
\noindent\textit{Source reference.} Original reference 0BB6: focused-seeds/geometry/S9-spaces-univ-catenary-dim-func/ref-tags.txt, source lines 7--7. Navigation pointer only: the counted proof cites this statement by \textbackslash{}ref as a standing input; the deduction the counted proof has to make is carried out inside the counted span, and no part of that deduction is delegated to this pointer.\par\medskip
\ifdefined\hypertarget\hypertarget{sourcepointer-5}{}\fi
\noindent\textit{Source reference.} Original reference 0ED5: focused-seeds/geometry/S9-spaces-univ-catenary-dim-func/ref-tags.txt, source lines 8--8. Navigation pointer only: the counted proof cites this statement by \textbackslash{}ref as a standing input; the deduction the counted proof has to make is carried out inside the counted span, and no part of that deduction is delegated to this pointer.\par\medskip
\ifdefined\hypertarget\hypertarget{sourcepointer-6}{}\fi
\noindent\textit{Source reference.} Original lemma 0EDM: sources/stacks-spaces-more-morphisms/pinned-original-source.tex, source lines 7274--7274. Navigation pointer only: the counted proof cites this statement by \textbackslash{}ref as a standing input; the deduction the counted proof has to make is carried out inside the counted span, and no part of that deduction is delegated to this pointer.\par\medskip
\ifdefined\hypertarget\hypertarget{sourcepointer-7}{}\fi
\noindent\textit{Source reference.} Original reference 02IL: focused-seeds/geometry/S9-spaces-univ-catenary-dim-func/ref-tags.txt, source lines 9--9. Navigation pointer only: the counted proof cites this statement by \textbackslash{}ref as a standing input; the deduction the counted proof has to make is carried out inside the counted span, and no part of that deduction is delegated to this pointer.\par\medskip
\ifdefined\hypertarget\hypertarget{sourcepointer-8}{}\fi
\noindent\textit{Source reference.} Original reference 09YC: focused-seeds/geometry/S9-spaces-univ-catenary-dim-func/ref-tags.txt, source lines 10--10. Navigation pointer only: the counted proof cites this statement by \textbackslash{}ref as a standing input; the deduction the counted proof has to make is carried out inside the counted span, and no part of that deduction is delegated to this pointer.\par\medskip
\ifdefined\hypertarget\hypertarget{sourcepointer-9}{}\fi
\noindent\textit{Source reference.} Original reference 04NM: focused-seeds/geometry/S9-spaces-univ-catenary-dim-func/ref-tags.txt, source lines 11--11. Navigation pointer only: the counted proof cites this statement by \textbackslash{}ref as a standing input; the deduction the counted proof has to make is carried out inside the counted span, and no part of that deduction is delegated to this pointer.\par\medskip
\ifdefined\hypertarget\hypertarget{sourcepointer-10}{}\fi
\noindent\textit{Source reference.} Original reference 04NR: focused-seeds/geometry/S9-spaces-univ-catenary-dim-func/ref-tags.txt, source lines 12--12. Navigation pointer only: the counted proof cites this statement by \textbackslash{}ref as a standing input; the deduction the counted proof has to make is carried out inside the counted span, and no part of that deduction is delegated to this pointer.\par\medskip
\ifdefined\hypertarget\hypertarget{sourcepointer-11}{}\fi
\noindent\textit{Source reference.} Original reference 04NZ: focused-seeds/geometry/S9-spaces-univ-catenary-dim-func/ref-tags.txt, source lines 13--13. Navigation pointer only: the counted proof cites this statement by \textbackslash{}ref as a standing input; the deduction the counted proof has to make is carried out inside the counted span, and no part of that deduction is delegated to this pointer.\par\medskip
\begin{lemma}
\label{lemma-universally-catenary-dimension-function}
Let $S$ be a scheme. Let $B$ be an algebraic space over $S$.
Let $\delta : |B| \to \mathbf{Z}$ be a function.
Assume $B$ is decent, locally Noetherian, and
universally catenary and $\delta$ is a dimension function.
If $X$ is a decent algebraic space over $B$ whose structure morphism
$f : X \to B$ is locally of finite type we define
$\delta_X : |X| \to \mathbf{Z}$ by the rule
$$
\delta_X(x) = \delta(f(x)) + \text{transcendence degreeof }x/f(x)
$$
(Morphisms of Spaces, Definition
\ref{spaces-morphisms-definition-dimension-fibre}).
Then $\delta_X$ is a dimension function.
\end{lemma}

\begin{proof}
The problem is local on $B$. Thus we may assume $B$ is quasi-compact.
By Decent Spaces, Lemma
\ref{decent-spaces-lemma-locally-Noetherian-decent-quasi-separated}
we see $B$ is quasi-separated. By
Limits of Spaces, Proposition
\ref{spaces-limits-proposition-there-is-a-scheme-finite-over}
we can choose a finite surjective morphism $\pi : Y \to X$
where $Y$ is a scheme.
Claim: $\delta_Y$ is a dimension function.

\medskip\noindent
The claim implies the lemma. With $X \to B$ as in the lemma
set $Z = Y \times_B X$ with projections $p : Z \to Y$ and $q : Z \to X$.
Then we have
$$
\delta_Z(z) = \delta_Y(p(z)) + \text{transcendence degreeof }z/p(z)
$$
and $\delta_Z(z) = \delta_X(q(z))$. This follows from
Morphisms of Spaces, Lemma 
\ref{spaces-morphisms-lemma-dimension-fibre-at-a-point-additive}
and the fact that these transcendence degrees are zero
for finite morphisms. By Decent Spaces, Lemma
\ref{decent-spaces-lemma-scheme-with-dimension-function}
and the claim we find that $\delta_Z$ is a dimension function.
Then we find that $\delta_X$ is a dimension function by
Decent Spaces, Lemma
\ref{decent-spaces-lemma-check-dimension-function-finite-cover}.

\medskip\noindent
Proof of the claim. Consider a specialization $y \leadsto y'$,
$y \not = y'$ of points of the Noetherian scheme $Y$.
Then $\delta_Y(y) > \delta_Y(y')$ because there are
no specializations between points in fibres of $Y$
(see Decent Spaces, Lemma
\ref{decent-spaces-lemma-conditions-on-fibre-and-qf}).
Using this for a chain of specializations we find
$$
\delta_Y(y) - \delta_Y(y') \geq
\text{codim}(\overline{\{y'\}}, \overline{\{y\}})
$$
Our task is to show equality. By
Properties, Lemma \ref{properties-lemma-locally-Noetherian-closed-point}
we can choose a specialization $y' \leadsto y_0$.
It suffices to show
$\delta_Y(y) - \delta_Y(y_0) =
\text{codim}(\overline{\{y_0\}}, \overline{\{y\}})$
because this will imply the equality for both
$y \leadsto y'$ and $y' \leadsto y_0$.

\medskip\noindent
Choose a maximal chain
$y = y_c \leadsto y_{c - 1} \leadsto \ldots \leadsto y_0$
of specializations in $Y$.
Set $b = \pi(y)$ and $b_0 = \pi(y_0)$.
Choose a maximal chain
$b = b_e \leadsto b_{e - 1} \leadsto \ldots \leadsto b_0$
of specializations in $|B|$.
We have to show $e = c$.
Since $\pi$ is closed
(Morphisms of Spaces, Lemma \ref{spaces-morphisms-lemma-finite-proper})
we can find a sequence of specializations
$y = y'_e \leadsto y'_{e - 1} \leadsto \ldots \leadsto y'_0$
mapping to
$b = b_e \leadsto b_{e - 1} \leadsto \ldots \leadsto b_0$.
Observe that $y'_e \leadsto y'_{e - 1} \leadsto \ldots \leadsto y'_0$
is a maximal chain as well.
If $y_0 = y'_0$, then because $Y$ is catenary, we conclude
that $e = c$ as desired. In the next paragraph we reduce to
this case by sleight of hand and we conclude in the same manner.

\medskip\noindent
Since $\pi$ is closed we see that $b_0$ is a closed point of $|B|$.
By Decent Spaces, Lemma \ref{decent-spaces-lemma-decent-space-closed-point}
we can represent $b_0$ by a closed immersion $b_0 : \Spec(k) \to B$.
By Lemma \ref{lemma-construct-glueing}
we can find a factorization
$$
Y \to Y' \to X
$$
with $\pi' : Y' \to X$ finite and $Y \to Y'$ a morphism which
map $y_0$ and $y'_0$ to the same point and is an isomorphism
away from the inverse image of $b_0$.
(Of course $Y'$ won't be a scheme but this doesn't matter for the
argument that follows.)
Clearly the maximal chains of specializations
$y_c \leadsto y_{c - 1} \leadsto \ldots \leadsto y_0$ and
$y'_e \leadsto y'_{e - 1} \leadsto \ldots \leadsto y'_0$
map to maximal chains of specializations in $Y'$ having
the same start and end.
Since $B$ is universally catenary, we see that
$|Y'|$ is catenary and we conclude as before.
\end{proof}

\end{document}
